\documentclass[UTF8,11pt]{ctexart}

\usepackage[a4paper,margin=1in]{geometry}
\usepackage{amsmath,amssymb,bm,mathtools}
\usepackage{booktabs,tabularx,array}
\usepackage{enumitem}
\usepackage[hidelinks]{hyperref}

\title{One-step Generative Transport via Endpoint-posterior Velocity Matching\\
\large 方法主线、已有基础与重点设计模块}
\author{梁凡}

\newcommand{\E}{\mathbb{E}}
\newcommand{\R}{\mathbb{R}}
\newcommand{\cN}{\mathcal{N}}
\newcommand{\cL}{\mathcal{L}}
\newcommand{\cH}{\mathcal{H}}
\newcommand{\dd}{\,\mathrm{d}}
\newcommand{\softmax}{\operatorname{softmax}}
\newcommand{\ESS}{\operatorname{ESS}}
\newcommand{\KL}{\operatorname{KL}}
\newcommand{\W}{\mathcal{W}}
\newcolumntype{Y}{>{\raggedright\arraybackslash}X}

\begin{document}

\maketitle

\section{一句话主线}

本文主线可以定为：
\begin{center}
\begin{minipage}{0.88\linewidth}
\centering
\textbf{One-step generative transport via endpoint-posterior velocity matching.}
\end{minipage}
\end{center}

中文表述是：\textbf{通过端点后验速度匹配来学习一步生成传输。}

它的关键不是再训练一个自由形式的速度网络 \(v_\theta(x,t)\)，而是要求速度场必须由一个低维到高维的一步生成器
\[
h_\theta:\R^d\to\R^D
\]
诱导出来。这样最终采样仍然是一步：
\[
u\sim\pi(u),\qquad x_{\mathrm{gen}}=h_\theta(u).
\]

因此，这条线与普通 flow matching 或 stochastic interpolants 的区别是：
\[
\text{ordinary velocity model: }v_\theta(x,t)
\]
被替换成
\[
\text{endpoint-induced velocity model: }
v_{\theta,\tau}(x)
=
\E_\theta[\dot X_\tau\mid X_\tau=x],
\]
其中条件期望的端点 \(X_1^\theta=h_\theta(U)\) 由同一个 one-step generator 给出。

\section{核心问题设定}

给定真实数据端点
\[
X_1\sim p_{\mathrm{data}},
\]
base endpoint
\[
X_0\sim p_0,
\]
以及 latent 变量
\[
U\sim\pi.
\]
我们希望学习
\[
X_1^\theta=h_\theta(U)
\]
使得 \(h_\theta(U)\) 接近 \(p_{\mathrm{data}}\)。

最简单的线性插值为
\[
X_\tau=(1-\tau)X_0+\tau X_1,\qquad \tau\in(0,1).
\]
真实 sample-wise velocity 是
\[
\dot X_\tau=X_1-X_0.
\]

模型端用一步生成器替代真实终点：
\[
X_{\theta,\tau}=(1-\tau)X_0+\tau h_\theta(U).
\]
模型速度场定义为
\[
v_{\theta,\tau}(x_\tau)
=
\E_\theta[h_\theta(U)-X_0\mid X_{\theta,\tau}=x_\tau].
\]

\section{端点后验形式}

在线性无噪声插值中，给定 \(x_\tau\) 和 \(U=u\)，可以反推出
\[
X_0
=
\frac{x_\tau-\tau h_\theta(u)}{1-\tau}.
\]
因此模型速度场可以写成
\begin{align}
v_{\theta,\tau}(x_\tau)
&=
\E_\theta
\left[
h_\theta(U)
-
\frac{x_\tau-\tau h_\theta(U)}{1-\tau}
\mid X_{\theta,\tau}=x_\tau
\right]\nonumber\\
&=
\frac{1}{1-\tau}
\E_\theta[h_\theta(U)\mid X_{\theta,\tau}=x_\tau]
-
\frac{1}{1-\tau}x_\tau.
\end{align}

令
\[
m_{\theta,\tau}(x_\tau)
:=
\E_\theta[h_\theta(U)\mid X_{\theta,\tau}=x_\tau],
\]
则
\[
v_{\theta,\tau}(x_\tau)
=
\frac{m_{\theta,\tau}(x_\tau)-x_\tau}{1-\tau}.
\]

真实 target 也可写为
\[
X_1-X_0
=
\frac{X_1-X_\tau}{1-\tau}.
\]
所以 velocity matching loss
\[
\cL_{\mathrm{vel}}(\theta)
=
\E
\left[
\left\|
v_{\theta,\tau}(X_\tau)-(X_1-X_0)
\right\|_2^2
\right]
\]
等价于
\[
\cL_{\mathrm{vel}}(\theta)
=
\frac{1}{(1-\tau)^2}
\E
\left[
\left\|
m_{\theta,\tau}(X_\tau)-X_1
\right\|_2^2
\right].
\]

\section{线性有噪声插值的具体推导}

上一节的推导依赖无噪声线性插值
\[
X_{\theta,\tau}=(1-\tau)X_0+\tau h_\theta(U).
\]
为了沿用第三节的记号，下文仍用 \(x_\tau\) 表示给定的中间层样本点；仅用上标 \(\gamma\)
区分带噪声的中间变量：
\[
X_{\theta,\tau}^{\gamma}
=
(1-\tau)X_0+\tau h_\theta(U)+\gamma_\tau\varepsilon,
\qquad
\varepsilon\sim\cN(0,I_D).
\]
其中 \(X_0\sim p_0\)、\(U\sim\pi\)，且三者相互独立。

\subsection{一般 joint posterior 形式}

给定 \(X_0=x_0\) 和 \(U=u\)，\(X_{\theta,\tau}^{\gamma}\) 的似然为
\[
L_\tau(x_\tau\mid x_0,u)
=
\phi\!\left(
x_\tau;\,
(1-\tau)x_0+\tau h_\theta(u),
\gamma_\tau^2 I_D
\right).
\]
因此 joint endpoint-posterior 满足
\[
p_\theta(x_0,u\mid x_\tau)
\propto
p_0(x_0)\pi(u)L_\tau(x_\tau\mid x_0,u).
\]

sample-wise velocity 为
\[
\dot X_{\theta,\tau}^{\gamma}
=
h_\theta(U)-X_0+\dot\gamma_\tau\varepsilon.
\]
由于在给定 \(x_\tau,x_0,u\) 时
\[
\varepsilon
=
\frac{
x_\tau-(1-\tau)x_0-\tau h_\theta(u)
}{
\gamma_\tau
},
\]
模型诱导的条件速度场可以写成
\begin{align}
v_{\theta,\tau}^{\gamma}(x_\tau)
&=
\E_\theta
\left[
h_\theta(U)-X_0+\dot\gamma_\tau\varepsilon
\mid X_{\theta,\tau}^{\gamma}=x_\tau
\right]\nonumber\\
&=
\E_\theta
\left[
h_\theta(U)-X_0
+\lambda_\tau
\left\{
x_\tau-(1-\tau)X_0-\tau h_\theta(U)
\right\}
\mid X_{\theta,\tau}^{\gamma}=x_\tau
\right],
\end{align}
其中
\[
\lambda_\tau:=\frac{\dot\gamma_\tau}{\gamma_\tau}.
\]

沿用第三节的 endpoint posterior mean 记号：
\[
m_{\theta,\tau}(x_\tau)
:=
\E_\theta[h_\theta(U)\mid X_{\theta,\tau}^{\gamma}=x_\tau].
\]
此外只额外定义一个 base endpoint posterior mean：
\[
m_{\theta,\tau}^{0}(x_\tau)
:=
\E_\theta[X_0\mid X_{\theta,\tau}^{\gamma}=x_\tau].
\]
则
\begin{align}
v_{\theta,\tau}^{\gamma}(x_\tau)
&=
m_{\theta,\tau}(x_\tau)-m_{\theta,\tau}^{0}(x_\tau)\nonumber\\
&\quad+
\lambda_\tau
\left[
x_\tau-(1-\tau)m_{\theta,\tau}^{0}(x_\tau)-\tau m_{\theta,\tau}(x_\tau)
\right].
\end{align}

同时，带噪声 bridge 的 score 为
\[
s_{\theta,\tau}^{\gamma}(x_\tau)
=
\nabla_{x_\tau}\log p_{\theta,\tau}^{\gamma}(x_\tau)
=
-
\frac{
x_\tau-(1-\tau)m_{\theta,\tau}^{0}(x_\tau)-\tau m_{\theta,\tau}(x_\tau)
}{
\gamma_\tau^2
}.
\]
因此速度场也可写成
\[
v_{\theta,\tau}^{\gamma}(x_\tau)
=
m_{\theta,\tau}(x_\tau)-m_{\theta,\tau}^{0}(x_\tau)
-\gamma_\tau\dot\gamma_\tau s_{\theta,\tau}^{\gamma}(x_\tau).
\]
这说明：有噪声线性插值下，velocity 不再只由第三节中的 endpoint posterior mean
\[
m_{\theta,\tau}(x_\tau)=
\E_\theta[h_\theta(U)\mid X_{\theta,\tau}^{\gamma}=x_\tau]
\]
决定，还会显式混入 base endpoint posterior mean \(m_{\theta,\tau}^{0}(x_\tau)\)
和 bridge score \(s_{\theta,\tau}^{\gamma}(x_\tau)\)。

若取
\[
\gamma_\tau=\sqrt{2\tau(1-\tau)},
\]
则
\[
\gamma_\tau\dot\gamma_\tau=1-2\tau,
\qquad
\lambda_\tau
=
\frac{1-2\tau}{2\tau(1-\tau)}.
\]
代回上式可得
\[
v_{\theta,\tau}^{\gamma}(x_\tau)
=
\frac{1}{2(1-\tau)}m_{\theta,\tau}(x_\tau)
-
\frac{1}{2\tau}m_{\theta,\tau}^{0}(x_\tau)
+
\frac{1-2\tau}{2\tau(1-\tau)}x_\tau.
\]
等价地，
\[
v_{\theta,\tau}^{\gamma}(x_\tau)
=
\frac{m_{\theta,\tau}(x_\tau)-x_\tau}{1-\tau}
-s_{\theta,\tau}^{\gamma}(x_\tau)
=
\frac{x_\tau-m_{\theta,\tau}^{0}(x_\tau)}{\tau}
+s_{\theta,\tau}^{\gamma}(x_\tau).
\]

\subsection{Gaussian base 下的低维 anchor 形式}

如果进一步取
\[
p_0=\cN(0,I_D),
\qquad
\gamma_\tau^2=2\tau(1-\tau),
\]
则可以把 \(X_0\) 积掉，只在 latent \(U\) 上做 anchors。
给定 \(U=u\) 时，
\[
X_{\theta,\tau}^{\gamma}\mid U=u
\sim
\cN
\left(
\tau h_\theta(u),
\left[(1-\tau)^2+\gamma_\tau^2\right]I_D
\right).
\]
记
\[
S_\tau=(1-\tau)^2+\gamma_\tau^2.
\]
则
\[
p_\theta(u\mid x_\tau)
\propto
\pi(u)
\phi
\left(
x_\tau;\tau h_\theta(u),S_\tau I_D
\right).
\]

在该线性高斯子问题中，
\[
\E[X_0\mid x_\tau,u]
=
\frac{1-\tau}{S_\tau}
\left(
x_\tau-\tau h_\theta(u)
\right),
\]
并且
\[
\E[\varepsilon\mid x_\tau,u]
=
\frac{\gamma_\tau}{S_\tau}
\left(
x_\tau-\tau h_\theta(u)
\right).
\]
所以给定 \(x_\tau,u\) 的 posterior velocity contribution 是
\[
a_\tau(x_\tau,u)
=
h_\theta(u)
+
\frac{
\gamma_\tau\dot\gamma_\tau-(1-\tau)
}{
S_\tau
}
\left(
x_\tau-\tau h_\theta(u)
\right).
\]

对特殊选择 \(\gamma_\tau^2=2\tau(1-\tau)\)，有
\[
S_\tau=1-\tau^2,
\qquad
\gamma_\tau\dot\gamma_\tau=1-2\tau.
\]
因此上式简化为
\[
a_\tau(x_\tau,u)
=
\frac{
h_\theta(u)-\tau x_\tau
}{
1-\tau^2
}.
\]
从而模型速度场为
\[
v_{\theta,\tau}^{\gamma}(x_\tau)
=
\frac{
m_{\theta,\tau}(x_\tau)-\tau x_\tau
}{
1-\tau^2
},
\]
其中
\[
m_{\theta,\tau}(x_\tau)
:=
\E_\theta[h_\theta(U)\mid X_{\theta,\tau}^{\gamma}=x_\tau].
\]

这给出一个很干净的 anchor 估计器。取 anchors \(u_k\sim q(u\mid x_\tau)\)，令
\[
h_k=h_\theta(u_k),
\]
则 logits 可取
\[
z_k
=
-
\frac{
\|x_\tau-\tau h_k\|_2^2
}{
2(1-\tau^2)
}
+
\log\pi(u_k)
-
\log q(u_k\mid x_\tau).
\]
归一化
\[
w_k=\frac{\exp(z_k)}{\sum_{j=1}^K\exp(z_j)}
\]
后，
\[
\widehat m_{\theta,K,\tau}(x_\tau)
=
\sum_{k=1}^K w_k h_\theta(u_k),
\qquad
\widehat v_{\theta,K,\tau}^{\gamma}(x_\tau)
=
\frac{
\widehat m_{\theta,K,\tau}(x_\tau)-\tau x_\tau
}{
1-\tau^2
}.
\]
若 \(q=\pi\)，则
\[
w_k
=
\softmax_k
\left(
-
\frac{
\|x_\tau-\tau h_\theta(u_k)\|_2^2
}{
2(1-\tau^2)
}
\right).
\]
和无噪声版本中的 \((1-\tau)^2\) 分母相比，这里的分母是
\(1-\tau^2\)。当 \(\tau\to 1\) 时，权重仍然会变尖，但尖锐程度从
\((1-\tau)^{-2}\) 降为 \((1-\tau)^{-1}\)，因此 anchors 的数值稳定性更好。

\subsection{对应训练目标}

真实数据端采用同样的 noisy bridge：
\[
X_\tau^{\gamma}
=
(1-\tau)X_0+\tau X_1+\gamma_\tau\varepsilon,
\]
sample-wise velocity 为
\[
V_\tau^{\gamma}
=
X_1-X_0+\dot\gamma_\tau\varepsilon.
\]
在 Gaussian base 和
\(\gamma_\tau^2=2\tau(1-\tau)\) 下，真实条件速度满足
\[
v_{\mathrm{data},\tau}^{\gamma}(x_\tau)
=
\frac{
m_{\mathrm{data},\tau}(x_\tau)-\tau x_\tau
}{
1-\tau^2
},
\qquad
m_{\mathrm{data},\tau}(x_\tau)
=
\E[X_1\mid X_\tau^{\gamma}=x_\tau].
\]

因此可训练
\[
\cL_{\mathrm{noisy}}(\theta)
=
\E
\left[
\left\|
\frac{
\widehat m_{\theta,K,\tau}(X_\tau^{\gamma})-\tau X_\tau^{\gamma}
}{
1-\tau^2
}
-V_\tau^{\gamma}
\right\|_2^2
\right].
\]
等价地，
\[
\cL_{\mathrm{noisy}}(\theta)
=
\frac{1}{(1-\tau^2)^2}
\E
\left[
\left\|
\widehat m_{\theta,K,\tau}(X_\tau^{\gamma})
-
\left\{
\tau X_\tau^{\gamma}+(1-\tau^2)V_\tau^{\gamma}
\right\}
\right\|_2^2
\right].
\]
其中随机目标
\[
T_\tau
:=
\tau X_\tau^{\gamma}+(1-\tau^2)V_\tau^{\gamma}
\]
满足
\[
\E[T_\tau\mid X_\tau^{\gamma}=x_\tau]
=
m_{\mathrm{data},\tau}(x_\tau).
\]
所以这一目标仍然是在学习 endpoint posterior mean，只是无噪声版本中的
\(\E[X_1\mid X_\tau=x]\) 被替换为带噪声 bridge 下的
\(\E[X_1\mid X_\tau^{\gamma}=x_\tau]\)。

\section{Anchor 近似}

后验密度满足
\[
p_\theta(u\mid x_\tau)
\propto
\pi(u)\,
p_0
\left(
\frac{x_\tau-\tau h_\theta(u)}{1-\tau}
\right).
\]

令
\[
h_k=h_\theta(u_k),
\qquad
x_{0,k}(x_\tau)
=
\frac{x_\tau-\tau h_k}{1-\tau}.
\]
importance weights 为
\[
\omega_k(x_\tau)
=
\pi(u_k)\,
p_0(x_{0,k}(x_\tau)).
\]

于是
\[
m_{\theta,K,\tau}(x_\tau)
=
\sum_{k=1}^K w_kh_\theta(u_k),
\]
并且
\[
v_{\theta,K,\tau}(x_\tau)
=
\frac{
\sum_{k=1}^K w_kh_\theta(u_k)-x_\tau
}{1-\tau}.
\]

当
\[
p_0=\cN(0,I_D),\qquad q(u\mid x_\tau)=\pi(u),
\]
权重化简为
\[
w_k
=
\softmax_k
\left(
-
\frac{
\|x_\tau-\tau h_\theta(u_k)\|_2^2
}{
2(1-\tau)^2
}
\right).
\]

直觉是：哪个 \(u_k\) 让反推出的
\[
x_{0,k}
=
\frac{x_\tau-\tau h_\theta(u_k)}{1-\tau}
\]
更像 base distribution \(p_0\) 中的样本，哪个 anchor 的权重就更大。


\section{重点设计一：Anchor proposal}

最基础的选择是
\[
u_k\sim\pi(u).
\]
但当 \(\tau\) 接近 1 时，权重项
\[
\exp\left(
-
\frac{\|x_\tau-\tau h_\theta(u_k)\|^2}{2(1-\tau)^2}
\right)
\]
非常尖锐，大多数 anchors 权重接近 0。

因此需要设计更强的 proposal。建议重点考虑三类。

\subsection{Amortized proposal}

训练一个条件 proposal：
\[
q_\phi(u\mid x_\tau).
\]
采样：
\[
u_k\sim q_\phi(u\mid x_\tau),
\]
权重仍然使用 importance correction：
\[
w_k\propto
\frac{
\pi(u_k)p_0(x_{0,k})
}{
q_\phi(u_k\mid x_\tau)
}.
\]

这可以用反向 KL 训练：
\[
\min_\phi
\E_{x_\tau}
\KL\left(
q_\phi(u\mid x_\tau)
\|p_\theta(u\mid x_\tau)
\right),
\]
或者用 self-normalized importance sampling 的有效样本数作为辅助目标。

\subsection{Mixture proposal}

为了防止 learned proposal 漏掉 mode，可以使用混合 proposal：
\[
q(u\mid x_\tau)
=
\lambda\pi(u)
+
(1-\lambda)q_\phi(u\mid x_\tau).
\]

prior 分量负责覆盖，learned proposal 负责精度。这个设计比纯 MAP-Laplace 更有研究空间，因为它可以处理多峰后验。

\subsection{Retrieval anchors}

维护一个 large codebook：
\[
\mathcal C=\{u_j,h_\theta(u_j)\}_{j=1}^M.
\]
给定 \(x_\tau\)，先检索使
\[
p_0\left(
\frac{x_\tau-\tau h_\theta(u_j)}{1-\tau}
\right)
\]
最大的 top-\(K\) anchors，再围绕它们做局部扰动。

这条路线工程上最容易先做出结果，也便于做 ESS 和速度的 ablation。

\section{重点设计二：\(\tau\) 选择}

\(\tau\) 的作用类似 MAGT 中的 smoothing level。过小的 \(\tau\) 使 \(X_\tau\) 太接近 base，训练 endpoint 的信号弱；过大的 \(\tau\) 使后验过尖，anchors 容易塌缩。

可以定义 effective sample size：
\[
\ESS(x_\tau)
=
\frac{1}{\sum_{k=1}^Kw_k(x_\tau)^2}.
\]

一个可设计的选择准则是：
\[
\tau^\star
=
\arg\min_{\tau\in\mathcal T}
\widehat{\W}_2
\left(
\{h_{\theta_\tau}(u_j)\},
\mathcal D_{\mathrm{val}}
\right)
\]
并加上约束
\[
\E[\ESS(X_\tau)]\ge \rho K.
\]

也可以把 \(\tau\) 做成 adaptive：
\[
\tau=\tau(x)
\]
或采用少量层级
\[
\tau\in\{\tau_1,\dots,\tau_L\},
\qquad L\ll 100.
\]
但如果目标是 one-step generation，建议第一版保持单一 \(\tau\)，把方法锋利地和 diffusion 区分开。

\section{重点设计三：Base endpoint}

若 \(X_0\sim\cN(0,I_D)\)，理论和实践都简单，但它是全维分布，可能削弱低维流形优势。

一个更贴近 MAGT 的设计是让 base endpoint 也由低维 latent 生成：
\[
X_0=g_\psi(V),
\qquad
V\sim p_V,
\qquad
g_\psi:\R^{d_0}\to\R^D.
\]
于是
\[
X_\tau=(1-\tau)g_\psi(V)+\tau h_\theta(U).
\]

如果进一步让 \(V=U\) 或让 \(V\) 与 \(U\) 相关，可以得到 shared-latent bridge：
\[
X_\tau=(1-\tau)g_\psi(U)+\tau h_\theta(U).
\]

这可能带来两个好处：
\begin{enumerate}
\item 中间路径仍主要受低维结构控制，有机会保留 intrinsic-dimension rate；
\item 条件后验从 \((X_0,U)\) 的全维搜索变成低维 latent 搜索，anchors 更高效。
\end{enumerate}

风险是 \(g_\psi\) 的引入会让模型更复杂，需要避免把问题变成另一个普通 autoencoder 或 GAN。

\section{重点设计四：带噪声 bridge 的取舍}

前面的具体推导说明，带噪声线性 bridge 并不只是给无噪声插值加一个扰动项。
一般情况下，速度场同时依赖
\[
m_{\theta,\tau}(x_\tau)=\E[h_\theta(U)\mid X_{\theta,\tau}^{\gamma}=x_\tau],
\qquad
m_{\theta,\tau}^{0}(x_\tau)=\E[X_0\mid X_{\theta,\tau}^{\gamma}=x_\tau],
\]
以及 bridge score \(s_{\theta,\tau}^{\gamma}(x_\tau)\)。这意味着如果直接采用一般 \(p_0\)，
anchors 可能需要在 \((x_0,u)\) 的 joint posterior 上工作。

但在
\[
p_0=\cN(0,I_D),
\qquad
\gamma_\tau^2=2\tau(1-\tau)
\]
这个特殊而自然的设定下，可以把 \(X_0\) 积掉，得到低维 latent posterior：
\[
p_\theta(u\mid x_\tau)
\propto
\pi(u)
\phi\left(x_\tau;\tau h_\theta(u),(1-\tau^2)I_D\right),
\]
并且速度场简化为
\[
v_{\theta,\tau}^{\gamma}(x_\tau)
=
\frac{m_{\theta,\tau}(x_\tau)-\tau x_\tau}{1-\tau^2}.
\]

因此带噪声版本有两个直接优势：
\begin{enumerate}
\item 权重分母从无噪声版本的 \((1-\tau)^2\) 变成 \(1-\tau^2\)，posterior anchors 在 \(\tau\to1\) 时更不容易完全塌缩；
\item 它保留了 endpoint posterior mean 的核心解释，同时更接近 stochastic interpolants 的连续 bridge 视角。
\end{enumerate}

主要风险是理论证明更难：无噪声版本可以直接反推 \(X_0\)，而带噪声版本需要处理 score 项或线性高斯消元带来的额外误差。
因此建议实验上优先比较无噪声 anchors 与 Gaussian-base noisy anchors，而理论上先证明后者的单层 conditional-mean 误差控制。

\section{可形成的论文贡献}

建议贡献写成以下三点。

\begin{enumerate}[label=\textbf{C\arabic*.}]
\item \textbf{Endpoint-posterior velocity matching.}
提出一种由 one-step generator 诱导速度场的训练目标，将 velocity matching 转化为 endpoint posterior matching。

\item \textbf{ESS-aware anchor approximation.}
提出面向 endpoint posterior 的 anchor proposal 设计，包括 amortized proposal、mixture proposal 或 retrieval anchors，并用 ESS 控制后验权重塌缩。

\item \textbf{Single-level velocity-to-generation analysis.}
证明或部分证明固定层级 velocity error 可以控制 endpoint generation error，并分析何时保留 intrinsic-dimension 优势。
\end{enumerate}

如果理论进度不够，C3 可以先弱化为：
\[
\text{training objective decomposition + finite-anchor consistency + empirical validation}.
\]

\section{实验设计}

实验应当围绕三个问题设计。

\subsection{问题一：是否能一步生成}

对比对象：
\[
\text{DDIM},\quad \text{flow matching},\quad \text{GAN},\quad \text{MAGT}.
\]
评价：
\[
W_2,\quad \text{FID},\quad \text{sampling time},\quad \text{NFE}.
\]

\subsection{问题二：是否更贴近流形}

在 synthetic manifolds 上报告：
\[
W_2,\qquad \text{off-manifold rate},\qquad \text{nearest-manifold distance}.
\]
这部分可以复用 MAGT 的 rings、spiral、moons、checker、helix、torus。

\subsection{问题三：anchor proposal 是否有效}

必须做 ablation：
\[
\text{prior anchors}
\quad
\text{QMC anchors}
\quad
\text{retrieval anchors}
\quad
\text{amortized proposal}
\quad
\text{mixture proposal}.
\]

报告：
\[
\ESS,\quad
\text{weight entropy},\quad
\text{validation }W_2,\quad
\text{training stability},\quad
\text{wall-clock cost}.
\]

\section{最小可行版本}

最小可行版本建议如下：
\begin{enumerate}
\item 使用无噪声线性插值：
\[
X_\tau=(1-\tau)X_0+\tau X_1.
\]
\item 使用 \(X_0\sim\cN(0,I_D)\) 和 latent anchors \(u_k\sim\pi\)。
\item 训练目标：
\[
\min_\theta
\E
\left[
\left\|
\sum_{k=1}^K w_k(X_\tau)h_\theta(u_k)
-
X_1
\right\|^2
\right].
\]
\item 增加 ESS 监控：
\[
\ESS(X_\tau)=1/\sum_kw_k^2.
\]
\item 做 retrieval anchors 或 mixture proposal 作为主要改进。
\end{enumerate}

这个版本足够验证核心命题：\textbf{是否可以只通过 endpoint-posterior velocity matching 学出 one-step generator。}

\section{推荐推进路线}

\begin{table}[h]
\centering
\small
\begin{tabularx}{\linewidth}{p{0.16\linewidth}p{0.35\linewidth}Y}
\toprule
阶段 & 目标 & 产出 \\
\midrule
Stage 1 & 复现基础 endpoint-posterior loss & synthetic manifolds 上跑通 one-step generation \\
Stage 2 & 加入 ESS-aware anchor proposal & 证明 ESS、稳定性、\(W_2\) 改善 \\
Stage 3 & 与 MAGT / FM / DDIM 对比 & 形成主要实验表格 \\
Stage 4 & 写理论分解 & velocity loss 正交分解、anchor consistency、初步 generalization bound \\
Stage 5 & 攻 single-level velocity pull-back & 形成理论亮点或作为未来工作 \\
\bottomrule
\end{tabularx}
\end{table}

\section{总结}

这条研究线最清晰的定位是：
\[
\boxed{
\text{用 endpoint posterior 约束速度场，用一步 transport 完成生成。}
}
\]

已有基础已经足够支撑方法雏形：MAGT 提供 one-step transport 和 anchors，stochastic interpolants 提供 velocity matching 语言，前面的推导给出 endpoint-posterior 等价式。

真正需要重点设计的是：
\begin{enumerate}
\item anchor proposal，尤其是 ESS-aware 或 amortized proposal；
\item \(\tau\) 选择机制；
\item base endpoint 是否保持低维结构；
\item 是否引入带噪声 bridge；
\item single-level velocity pull-back 理论。
\end{enumerate}

其中最适合作为论文创新点的是：
\begin{center}
\fbox{
\begin{minipage}{0.82\linewidth}
\centering
endpoint-posterior velocity matching\\
\(+\) ESS-aware anchor proposal\\
\(+\) single-level velocity analysis
\end{minipage}
}
\end{center}

\end{document}
